% =====================================================================
% VERSION ENSAE, REDUCTION SFE (phase 3, 1er octobre 2026) : chapitre 6, modelisation multi-etats de la conformite.
% Copie a texte constant de chapitres/11_conformite_multietats.tex.
% Ne changent que les titres et les renvois repares, listes dans
% reduction/phase3.md.
% Lu par main_ensae.tex seul : ne jamais le faire lire par main_institut.tex,
% qui garde chapitres/.
% =====================================================================

\chapter{Modélisation multi-états de la conformité}
\label{chap:multietats}


\section{Deux échelles de temps}

Le besoin de capital dépend de l'état de conformité de l'entité, qui est gradué et évolue dans le
temps. Nous le modélisons à deux échelles. La couche rapide est celle des incidents et de leur
cascade dans l'année, chiffrée par le moteur du chapitre~\ref{chap:cascade} conditionnellement à
une configuration de conformité. Un Vasicek multi-états y classe chaque pilier comme non
conforme (NC), partiellement conforme (PC) ou conforme (C). La couche lente est celle de la
gouvernance, où une chaîne de Markov fait progresser la conformité sur des mois et des années, et son
état sert de paramètre à la couche rapide.

La structure est celle d'un modèle de crédit à migration de rating (\S\ref{sec:couplage}). Nous
séparons les deux horloges afin d'empêcher qu'une tendance longue explique un enchaînement de
quelques heures. Le préprint de cascade climatique \citep{ccrn2026} distingue aussi
le temps calendaire du pas de propagation.

\section{Vasicek multi-états à seuils ordonnés}

Chaque pilier possède une capacité de conformité latente $C^\star\sim\mathcal N(0,1)$, que nous
décomposons à la Vasicek en un facteur systémique commun $\Theta$ et un choc idiosyncratique :
\begin{equation}
C^\star \;=\; \gamma\,\Theta \;+\; \sqrt{1-\gamma^2}\,\varepsilon, \qquad \gamma=0{,}68.
\end{equation}
Au lieu d'un unique seuil de défaut, nous en posons deux, ordonnés, qui découpent trois
états :
\begin{equation}
\text{NC si } C^\star\le K_{\text{bas}}, \quad
\text{PC si } K_{\text{bas}}<C^\star\le K_{\text{haut}}, \quad
\text{C si } C^\star>K_{\text{haut}}.
\end{equation}
La charge $\gamma$ et les probabilités d'état sont posées à partir de proxies externes
(annexe~B). Pour NC, PC et C, nous retenons $0{,}35$, $0{,}35$ et $0{,}30$, d'où $K_{\text{bas}}=\Phi^{-1}(p_{\text{NC}})$ et
$K_{\text{haut}}=\Phi^{-1}(p_{\text{NC}}+p_{\text{PC}})$. Sachant l'environnement $\Theta$, on
obtient
\begin{equation}
\mathbb P(\text{NC}\mid\Theta)=\Phi\!\Big(\tfrac{K_{\text{bas}}-\gamma\Theta}{\sqrt{1-\gamma^2}}\Big),
\qquad
\mathbb P(\text{C}\mid\Theta)=1-\Phi\!\Big(\tfrac{K_{\text{haut}}-\gamma\Theta}{\sqrt{1-\gamma^2}}\Big).
\end{equation}
Intégrée sur $\Theta$, la marginale redonne l'ancrage, tandis qu'en crise systémique la masse bascule
vers la non-conformité (figure~\ref{fig:multietats-rapport}).

\begin{figure}[H]\centering
\figover{S10_scr_multi_etats.png}
\caption{États de conformité et capital. (A)~Besoin de capital du périmètre financier étudié par état, en cumulant la fréquence seule,
puis la propagation, puis la détection, trois des canaux définis au
\S\ref{sec:quatre-canaux}. (B)~Probabilités d'état en régime normal et en crise systémique : la
masse bascule vers la non-conformité.}
% SOURCES-SCRIPTS: 16
\label{fig:multietats-rapport}
\end{figure}

% SOURCES-SCRIPTS: 63


\section{Quatre canaux de la conformité}
\label{sec:quatre-canaux}

L'état de conformité agit sur la cascade par quatre canaux, l'agrégation restant inchangée. Le
tableau donne la valeur de chacun aux deux états extrêmes, où tous les piliers partagent le même état
(\S\ref{sec:etat-par-pilier}).

\begin{center}\footnotesize
\renewcommand{\arraystretch}{1.2}
\begin{tabular}{@{}l l r r >{\raggedright\arraybackslash}p{4.6cm}@{}}
\toprule
\textbf{Canal} & \textbf{Paramètre} & \textbf{Conforme} & \textbf{Non conforme} & \textbf{Statut} \\
\midrule
Fréquence & $\lambda$, incidents enregistrés par an & $21{,}56$ & $53{,}62$ & multiplicateur posé \\
Propagation & $g$, gain de la normalisation & $0{,}45$ & $0{,}90$ & valeurs posées \\
Détection & $p_u$, dépassement par pilier touché & $0{,}128$ & $0{,}181$ & multiplicateurs $0{,}85$ et $1{,}20$ posés \\
Accumulation & probabilité de choc commun & $0$ & $0{,}68$ & borne haute posée, non identifiée \\
\bottomrule
\end{tabular}
\end{center}

Chacun des quatre est décrit ci-dessous.
\begin{enumerate}[leftmargin=1.5em,itemsep=2pt,topsep=2pt]
\item L'état module la \emph{fréquence} $\lambda$ des incidents amorcés par un multiplicateur de
  scénario, de $21{,}56$ incidents enregistrés par an à l'état conforme à $53{,}62$ à l'état non
  conforme.
\item Le gain de \emph{propagation} $g$ de la normalisation~\eqref{eq:norm} baisse de $0{,}90$ à
  $0{,}45$, car une entité plus conforme contient mieux la contagion. La propension du pilier $j$ à
  propager vaut $e_j=g(c_j)\,s_j/\max_k s_k$, où $s_j$ est la somme de la ligne $j$ de TRANS et
  $c_j$ l'état du pilier $j$.
\item La \emph{détection} agit sur $p_u$, le taux de dépassement, c'est-à-dire la probabilité
  qu'une sévérité élémentaire dépasse le seuil $u$ du chapitre~\ref{chap:socle}. Le moteur l'applique à chaque pilier touché, si bien que la probabilité de perte matérielle d'un incident est plus élevée et dépend aussi de la propagation. Le taux $p_u$ varie avec l'état parce qu'une entité conforme intercepte et remédie tôt, de sorte qu'une
  moindre fraction des piliers touchés atteint la sévérité matérielle. Nous multiplions $p_u$ par $0{,}85$
  à l'état conforme et par $1{,}20$ à l'état non conforme, deux valeurs posées.
\item L'\emph{accumulation par un prestataire commun} transforme, à l'état non conforme, un
  incident amorcé sur le pilier des tiers en choc commun. Ce choc touche chacun des autres piliers
  avec la probabilité $0{,}68$, borne haute posée égale à la charge $\gamma$, cette probabilité
  n'étant pas identifiée (\S\ref{sec:p4-accumulation}).
\end{enumerate}
Le résultat principal du chapitre~\ref{chap:resultats} fait varier les quatre canaux entre les deux états. Le multiplicateur de
fréquence se compose vecteur par vecteur selon $\lambda_{\text{NC}}=\lambda_{\text{C}}\sum_v \pi_v\,m_v$, où $\pi_v$ est la
part du vecteur $v$ dans le recensement Hackmageddon (\S\ref{sec:hackmageddon}) et $m_v$ le centre
d'une fourchette de multiplicateur propre au scénario non conforme.

\begin{center}\footnotesize
\renewcommand{\arraystretch}{1.2}
\begin{tabular}{@{}l r c r r >{\raggedright\arraybackslash}p{5.2cm}@{}}
\toprule
\textbf{Vecteur} & \textbf{Part $\pi_v$} & \textbf{Fourchette} & \textbf{Centre} & \textbf{$\pi_v\,m_v$} & \textbf{Effet publié qui ancre la fourchette} \\
\midrule
Hameçonnage & $0{,}388$ & $[1{,}8\,;3{,}0]$ & $2{,}40$ & $0{,}9312$ & réduction du risque par la formation (Ponemon/IBM) \\
Exploitation de vulnérabilité & $0{,}338$ & $[2{,}0\,;2{,}6]$ & $2{,}30$ & $0{,}7774$ & écart de conversion en intrusion entre exploitation et hameçonnage (ENISA) \\
Chaîne d'approvisionnement & $0{,}158$ & $[1{,}9\,;2{,}6]$ & $2{,}25$ & $0{,}3555$ & surcoût et progression des brèches d'origine tierce (IBM, SecurityScorecard) \\
Identifiants compromis & $0{,}063$ & $[3{,}0\,;8{,}0]$ & $5{,}50$ & $0{,}3465$ & réduction de compromission par l'authentification forte (Microsoft Research) \\
Autres & $0{,}053$ & $[1{,}3\,;1{,}6]$ & $1{,}45$ & $0{,}0769$ & aucune source dédiée \\
\midrule
\multicolumn{4}{@{}l}{Multiplicateur agrégé, état non conforme} & $2{,}4875$ & $21{,}56\times 2{,}4875 = 53{,}62$ incidents par an \\
\bottomrule
\end{tabular}
\end{center}

\noindent Ce multiplicateur est posé, non calibré : les parts sont une citation externe, et
chaque fourchette traduit un effet publié en multiplicateur prudent. Aux bornes des fourchettes, il
va de $1{,}933$ à $3{,}042$, ce qui donne une fréquence non conforme de $41{,}7$ à $65{,}6$ incidents
par an. Le vecteur des identifiants, avec $6{,}3\,\%$ des incidents, représente $13{,}9\,\%$ du
multiplicateur.

% SOURCES-SCRIPTS: 67 99 43


\section{Valeurs posées de \texorpdfstring{$g$}{g} et invariance de l'ordre}
\label{sec:invariance-g}

Les trois valeurs du gain de propagation,
\[
  g = 0{,}45 \ \text{(conforme)}, \qquad 0{,}68 \ \text{(partiellement conforme)},
  \qquad 0{,}90 \ \text{(non conforme)},
\]
sont posées : aucune donnée ne donne la relation état $\to g$, les états DORA n'étant pas publiés.
Pour une entité de $15\,000$~M\euro{} de provisions (\S\ref{sec:ancrage}), le besoin passe de $116$ à $169$~M\euro{}
quand $g$ passe de $0{,}45$ à $0{,}90$. Le montant de $169$ reste une extrapolation d'ordre de grandeur.

La croissance du capital en $g$ est contrôlée numériquement.
Sur une grille de vingt points, à nombres aléatoires communs, aucun pas n'est décroissant, pour le
capital comme pour la perte moyenne, à l'échelle d'entité comme du secteur. Les applications testées respectant $g_C < g_{PC} < g_{NC}$ donnent $\mathrm{SCR}_C < \mathrm{SCR}_{PC} < \mathrm{SCR}_{NC}$. La grille ne démontre pas cette propriété pour toute valeur de $g$. L'existence et le sens de l'écart sont
ainsi vérifiés numériquement sur ces vingt points et sur les cinq applications du tableau, et seule
son amplitude est un scénario.

\begin{center}\small
\renewcommand{\arraystretch}{1.25}
\begin{tabular}{@{}l r r r r r@{}}
\toprule
\textbf{Application état $\to g$} & \textbf{Conf.} & \textbf{Part.} & \textbf{Non conf.} &
\textbf{Écart} & \textbf{Ordre} \\
 & M\euro & M\euro & M\euro & & \\
\midrule
étroite $(0{,}70 / 0{,}80 / 0{,}90)$ & $145{,}5$ & $156{,}5$ & $169{,}0$ & $+16{,}2\,\%$ & oui \\
\textbf{retenue} $(0{,}45 / 0{,}68 / 0{,}90)$ & $\mathbf{116{,}4}$ & $\mathbf{143{,}2}$ &
$\mathbf{169{,}0}$ & $\mathbf{+45{,}2\,\%}$ & oui \\
large $(0{,}20 / 0{,}55 / 0{,}90)$ & $94{,}4$ & $126{,}4$ & $169{,}0$ & $+79{,}1\,\%$ & oui \\
conformité faible $(0{,}45 / 0{,}55 / 0{,}65)$ & $116{,}4$ & $126{,}4$ & $139{,}2$ &
$+19{,}5\,\%$ & oui \\
très étroite $(0{,}85 / 0{,}875 / 0{,}90)$ & $162{,}6$ & $166{,}7$ & $169{,}0$ & $+3{,}9\,\%$ &
oui \\
\midrule
\emph{repère forfaitaire de Formule Standard} & $450$ & $450$ & $450$ & $0\,\%$ & \emph{plat} \\
\bottomrule
\end{tabular}
\end{center}

\vspace{2pt}
\noindent Avec un déplacement de $0{,}05$ sur $g$, l'écart est de $+3{,}9\,\%$, le forfait restant plat. L'hypothèse qui demeure est qualitative, $g_C < g_{NC}$, et signifie que se conformer à
DORA réduit la propagation entre piliers. L'élasticité d'arc du capital au gain est de $0{,}41$ sur
$g \in [0{,}20\,;0{,}95]$, et un écart de $10\,\%$ ne demande qu'une baisse de $g$ de $0{,}90$ à
$0{,}75$. Nous présentons donc l'écart entre états comme un scénario sur l'écart entre les valeurs de $g$ au
chapitre~\ref{chap:resultats}.
% SOURCES-SCRIPTS: 66 67

\section{État de conformité par pilier}
\label{sec:etat-par-pilier}

Chaque pilier reçoit sa propre latente $C^\star_j$, les cinq étant corrélées par le facteur commun
$\Theta$. Un pilier non conforme propage alors plus et détecte moins que ses voisins conformes. Les
canaux deviennent $\lambda_j$, $g_j$ et $p_{u,j}$, indexés par pilier. Le moteur redonne
la version globale lorsque tous les piliers partagent un état, et basculer un pilier à la
fois permet de décomposer le surcoût de non-conformité par pilier (chapitre~\ref{chap:resultats}).

\subsection{Dépendance entre états}
\label{sec:dependance-etats}

% SOURCES-SCRIPTS: 69 77 94

La conformité aux piliers~1 et~2 pourrait aider à satisfaire les exigences du pilier~3 si les contrôles se recouvrent. Cette hypothèse n'est pas observée. Les
cinq latentes partagent le facteur systémique avec une charge de $0{,}68$, soit une corrélation de
$0{,}462$ entre deux piliers quelconques. Cette dépendance est identique pour toutes les paires.

Ajouter un surcroît
de corrélation $\delta$ aux deux seules paires concernées, à marges fixées, est admissible jusqu'à
$\delta = 0{,}38$. Sur ce balayage, la probabilité que les cinq piliers soient non conformes passe
de $0{,}068$ à $0{,}088$ et celle qu'aucun ne soit non conforme de $0{,}291$ à $0{,}330$. Ce diagnostic est binaire : son état C regroupe PC et C, contrairement au modèle à trois états. Le capital
espéré, pour le périmètre financier étudié, ne varie que de $2$~M\euro, soit $0{,}02\,\%$.

Dans les $32$ configurations, la fréquence, la propagation et la détection varient, tandis que l'accumulation reste fixe.
Les capitaux vont de $6\,015$ à $17\,009$~M\euro. Une forme additive en indicateurs de pilier les approche avec $R^2 = 0{,}9945$.
Les contributions sont de $3\,524$ (P1), $2\,863$ (P4),
$2\,022$ (P2), $1\,577$ (P3) et $897$~M\euro{} (P5) au-dessus d'un socle tous conformes.
L'espérance d'une fonction additive ne dépend que des marges, donc l'invariance tient pour toute
structure de dépendance à marges fixées. Cette quasi-additivité est le même fait que la faible
interaction entre piliers, de signe non résolu (section sur la décomposition par pilier,
chapitre~\ref{chap:resultats}).

L'argument concerne une
espérance et ne se transpose ni au quantile de la loi de perte mélangée, où l'invariance est
seulement non détectée, ni au quantile de la loi des configurations, où l'effet est de quelques pour
cent.

\section{Chaîne de Markov et séjours de type phase}

% SOURCES-SCRIPTS: 17 77 08b

À l'échelle lente, l'état suit une chaîne de Markov à temps continu
$\text{NC}\to\text{PC}\to\text{C}$, dont l'état conforme est absorbant dans le cas de base. Le
retour vers un état moins conforme est traité en sensibilité. Nous posons des séjours d'Erlang à $k=2$ phases,
de durée moyenne $1{,}5$ an par état transitoire. Le taux par phase est $\mu_0=k/(1{,}5\ \text{an})$ dans le régime de référence.
Cette loi réduit les séjours très courts sans imposer de durée minimale. On développe chaque
macro-état transitoire en $k$ phases, et l'on résout le générateur $Q$ par exponentielle de matrice :
\begin{equation}
\big(\mathbb P(\text{NC},t),\mathbb P(\text{PC},t),\mathbb P(\text{C},t)\big)
\;=\; \text{(sommes de phases de)}\;\; v_0\,\exp(Qt),
\end{equation}
où $v_0$ place les masses $0{,}35$, $0{,}35$ et $0{,}30$ en première phase de NC, PC et en C. Cette distribution pondère les besoins de capital par état :
\begin{equation}
\overline{\mathrm{SCR}}(t)\;=\;\sum_{e\in\{\text{NC,PC,C}\}}\mathbb P(e,t)\,\mathrm{SCR}_e .
\end{equation}
Il s'agit d'une moyenne de quantiles, qui diffère du quantile à $99{,}5\,\%$ de la loi mélangée,
dont la queue vient des composantes les plus lourdes. Sur les $32$ configurations, le quantile du mélange dépasse la moyenne des quantiles de $9{,}3\,\%$. La trajectoire se lit donc comme l'espérance du
besoin de capital sous l'incertitude d'état.

% SOURCES-SCRIPTS: 08b 77

\section{Couplage systémique et analogie de crédit}
\label{sec:couplage}

Le même facteur $\Theta$ agit sur les deux couches. Un environnement dégradé, où $\Theta<0$,
ralentit la remédiation par des taux de transition $\mu(\Theta)=\mu_0\,e^{\beta_{\text{rem}}\Theta}$
et durcit la cascade par $\mathrm{SCR}_e(\Theta)=\mathrm{SCR}_e\,e^{-\beta_{\text{casc}}\Theta}$,
deux formes exponentielles posées, avec $\beta_{\text{rem}}=0{,}35$ et $\beta_{\text{casc}}=0{,}1$. Elles introduisent la corrélation entre conformité et
sinistralité, et la dispersion de $\Theta$
engendre la bande d'incertitude de la trajectoire. La chaîne de Markov joue le rôle de la migration
d'un rating de conformité, et le Vasicek multi-états celui du défaut conditionnel au rating. Les
résultats chiffrés figurent au chapitre~\ref{chap:resultats}.

% SOURCES-SCRIPTS: 63 17
